% Faithful official CCC2022 article17 mathematical transcription: necessary Section2 and full Section3 main proof.
\section*{2 Preliminaries}
\subsection*{2.1 Nisan–Wigderson generators}
For the rest of the paper for any two real numbers $r_1\leq r_2$, define $[r_1,r_2):=\{i\in\mathbb N:\lfloor r_1\rfloor\leq i<\lceil r_2\rceil\}$ and $[r_1,r_2]:=\{i\in\mathbb N:\lfloor r_1\rfloor\leq i\leq\lceil r_2\rceil\}$.
Let $f:\{0,1\}^*\to\{0,1\}$ be a Boolean function. For a natural number $n$, $f_n$ denotes the function $f$ restricted to $\{0,1\}^n$. Let $A$ be an $m\times n$ $0$-$1$ matrix such that each row of $A$ has exactly $l$ ones. Such a matrix is called an $l$-sparse matrix. For such a $m\times n$ $l$-sparse matrix $A$, $J_i(A):=\{j\in[0,n):A_{i,j}=1\}$.
For every pair $(f,A)$ where $f:\{0,1\}^l\to\{0,1\}$ is a Boolean function and $A$ is a $m\times n$ $l$-sparse matrix, Nisan and Wigderson [26] defined the generator $\mathrm{NW}_{f,A}:\{0,1\}^n\to\{0,1\}^m$ as follows:
For every input $a\in\{0,1\}^n$, the $i$'th bit of the output of $\mathrm{NW}_{f,A}(a)$ is $f(a|_{J_i(A)})$.
It was proved in the seminal paper [26] that if $f$ is a hard function (depending on the application) and $A$ satisfies specific combinatorial properties, then $\mathrm{NW}_{f,A}$ is a good pseudorandom generator (depending on the parameters).
Let $f\in\mathrm{NP}\cap\mathrm{CoNP}$. A pair of propositional formulas $(\sigma_0(p,q),\sigma_1(p,r))$ is a representation of $f_n$ for a natural number $n$ iff:
\begin{enumerate}
\item $|p|=n$ and moreover $p,q$, and $r$ variables are disjoint.
\item $(\sigma_0,\sigma_1)$ defines the function $f_n$ which means:
\begin{enumerate}
\item $\neg\sigma_0\vee\neg\sigma_1$ is a tautology.
\item For every $a\in\{0,1\}^n$, $f(a)=i$ iff $\sigma_i(a,t)$ is satisfiable where $i\in\{0,1\}$.
\end{enumerate}
\end{enumerate}
Note that as $f\in\mathrm{NP}\cap\mathrm{CoNP}$, for every $n$, $f_n$ has a representation.
Suppose $f\in\mathrm{NP}\cap\mathrm{CoNP}$, $(\sigma_0,\sigma_1)$ is a representation for $f_l$, and $A$ is a $m\times n$ $l$-sparse matrix. Then for any $b\in\{0,1\}^m$, $\tau_b(\mathrm{NW}_{f,A})$ based on $(\sigma_0,\sigma_1)$ is the following propositional formula:
\[\bigvee_{b_i=1}\neg\sigma_1(p|_{J_i(A)},q_i)\vee\bigvee_{b_i=0}\neg\sigma_0(p|_{J_i(A)},q_i)\]
where $q_i$'s are disjoint variables. Note that if $b\notin\operatorname{Rng}(\mathrm{NW}_{f,A})$, then $\tau_b(\mathrm{NW}_{f,A})$ is tautology.
As it was discussed in previous works [2,15,37,19], $\mathrm{NW}_{f,A}$ can be a hard proof complexity generator for a proof system $P$ for the following four reasons:
\begin{itemize}
\item The complexity of $f$.
\item The properties that $A$ satisfies.
\item The representation of $f$ that is used in the formula $\tau_b(\mathrm{NW}_{f,A})$.
\item The string $b\notin\operatorname{Rng}(\mathrm{NW}_{f,A})$.
\end{itemize}
As we will see, our main result also imposes different conditions on $\tau_b(\mathrm{NW}_{f,A})$ to make sure that it requires long proofs.
The following parts explain the properties that we need for the matrices and representations to prove our theorems.
\subsubsection*{2.1.1 Representations}
The hardness of $\tau_b(\mathrm{NW}_{f,A})$ can depend on the pair $(\sigma_0,\sigma_1)$ that is used in it. This matter has been investigated in [2], and they examined different representations. Recently, Sokolov [38] answered one of the open problems that was stated about a representation of $\tau_b(\mathrm{NW}_{f,A})$ in [2]. Here we investigate representations based on definability over finite structures in logic, which is a well-studied concept in descriptive complexity and finite model theory.
\paragraph{$\Sigma_1^1\cap\Pi_1^1$ representation.} Let $\mathcal L$ be a finite relational language and $X$ be a unary relational symbol which is not in $\mathcal L$. A $\Sigma_1^1$ formula $\psi(X)$ in the language $\mathcal L\cup\{X\}$ with equality defines a function $f\in\mathrm{NP}$ iff:
\begin{enumerate}
\item $\psi:=\exists\overline Y\ \varphi(X,\overline Y)$ where $\varphi(X,\overline Y)$ is a first-order formula in the language $\mathcal L\cup\{X\}$ with equality.
\item $X$ is not in $\overline Y$.
\item For every $n$, every $a\in\{0,1\}^n$, $f_n(a)=1$ iff $([0,n),a)\models\psi(X)$ when $X$ is interpreted by $a$.
\end{enumerate}
Fagin's theorem [6] directly implies that for every symmetric $f\in\mathrm{NP}$, a $\Sigma_1^1$ formula $\psi_f(X)$ exists in a language $\mathcal L\cup\{X\}$ that defines $f$. Therefore, the set of functions that are $\Sigma_1^1$ definable is exactly symmetric NP and hence this set is quite rich.
To prove Theorem 2, the following lemma is needed. This lemma states that the truth of first-order formulas in a relational language does not change under permutations.
If $A$ is a set and $Q$ is a relation on it, i.e. $Q\subseteq A^k$ for some $k$, and $h:A\to A$ is a function, then $h(Q):=\{(h(a_0),\ldots,h(a_{k-1})):(a_0,\ldots,a_{k-1})\in Q\}$.
\paragraph{Lemma 3.} Let $\mathcal L=\{Y_0,\ldots,Y_k\}$ be a finite relational language and $\mathcal A_0=(A,\{Q_0^0,\ldots,Q_0^k\})$ be an $\mathcal L$-structure. Let $h$ be a bijective function from $A$ onto $A$. Consider the $\mathcal L$-structure $\mathcal A_1:=(A,\{Q_1^0,\ldots,Q_1^k\})$ where $Q_1^i=h(Q_0^i)$, for every $i\in[0,k]$. Then for every first-order formula $\varphi(x_0,\ldots,x_{p-1})$ in $\mathcal L$ with equality, every $(a_0,\ldots,a_{p-1})\in A^p$:
\[\mathcal A_0\models\varphi(a_0,\ldots,a_{p-1})\Longleftrightarrow\mathcal A_1\models\varphi(h(a_0),\ldots,h(a_{p-1})).\]
\paragraph{Proof.} This lemma can be proved by induction on the complexity of $\varphi$.\hfill$\triangleleft$
Let $\exists\overline Y\ \varphi(X,\overline Y)$ be a $\Sigma_1^1$ formula. Then for any $n$, the Paris–Wilkie translation [27] (see also Section 8.2 of [21]) of $\varphi^{<n}(X,\overline Y)$ ($\varphi^{<n}$ is $\varphi$ when every first-order quantifier is bounded by $n$) is denoted by $\langle\varphi\rangle_n(p,q)$ which is a constant depth formula (without loss of generality we can assume that it is a CNF using extension variables). The number $n$ indicates the size of the universe in which $\varphi(X,Y)$ has been considered.
Let $f\in\mathrm{NP}\cap\mathrm{CoNP}$ be a symmetric function. Then a pair of $\Sigma_1^1$ formulas $(\exists\overline Y\ \varphi_0(X,\overline Y),\exists\overline Z\ \varphi_1(X,\overline Z))$ defines $f$ iff:
\begin{enumerate}
\item $\exists\overline Y\ \varphi_1(X,\overline Y)$ defines $f$.
\item $\exists\overline Z\ \varphi_0(X,\overline Z)$ defines $\neg f$.
\end{enumerate}
Such a pair is called a $\Sigma_1^1\cap\Pi_1^1$ definition of $f$. Moreover, for any $n$, $(\langle\varphi_0\rangle_n,\langle\varphi_1\rangle_n)$ is a representation of $f_n$. For the sake of easiness, by $\langle\psi\rangle_n$ we mean $\langle\varphi\rangle_n$ where $\psi(X):=\exists\overline Y\ \varphi(X,\overline Y)$ is a $\Sigma_1^1$ formula.
\subsection*{2.2 Proof systems}
We assume the reader knows the basic facts about proof complexity, proof systems, and bounded arithmetics (for a detailed discussion, see [21,9]). Here we state some useful facts about $\mathrm{AC}^0$-Frege, which will be used in the results.
\subsubsection*{2.2.1 $\mathrm{AC}^0$-Frege}
$\mathrm{AC}^0$-Frege is the name for a family of proof systems that work with constant-depth de Morgan formulas. For each $d\geq1$, $F_d$ denotes $\mathrm{AC}^0$-Frege proof system of depth $d$, which is the Frege proof system that works with formulas of depth at most $d$.
To prove Theorem 2, we need some known relations between $\mathrm{AC}^0$-Frege and $V_1^0$, which is a two-sorted bounded arithmetic (see [5,9]). These relations are related to the model theory of $V_1^0$.
Let $M$ be an arbitrary nonstandard model of true arithmetic and $n\in M\setminus\mathbb N$. Then
\[M_n:=\{a\in M:\text{There exists a }b\in M\setminus\mathbb N\text{ such that }a<2^{n^{1/b}}\}.\]
The following theorems explain the relationship between $\mathrm{AC}^0$-Frege and $V_1^0$ from the point of view of proof complexity. These theorems state that lower bounds for $\mathrm{AC}^0$-Frege correspond to unprovability in $V_1^0$ in a certain sense.
For a set $A$, $\mathcal P(A)$ denotes the power set of $A$.
\paragraph{Theorem 4 (Section 9.4 of [9]).} Let $(M,\chi)\models V_1^0$ and $\sigma\in\chi$ be a constant depth propositional formula (depth of $\sigma$ is standard). If $\neg\sigma$ is satisfiable by an assignment in $\chi$, then for every standard $d$, there is no $F_d$-proof of $\sigma$ in $(M,\chi)$.
Note that Theorem 4 also holds in the case where $\sigma$ is the Paris-Wilkie translation of a bounded arithmetical formula such as $\varphi(x,\overline R)$ ($\sigma=\langle\varphi(n,\overline R)\rangle_n$ for some $n\in M$), i.e. if there is an $\overline\alpha\in\chi$ such that $(M,\chi)\models\neg\varphi(n,\overline\alpha)$, then $\neg\sigma$ is satisfiable by an assignment from $\chi$ and therefore $\sigma$ does not have any $F_d$-proof in $M$.
\paragraph{Theorem 5 (Section 9.4 of [9]).} Let $M$ be a countable nonstandard model of true arithmetic and $\varphi(x,R)$ be a bounded arithmetical formula such that for every $d$, the family $\{\langle\varphi(n,R)\rangle_n\}_{n\in\mathbb N}$ requires exponential $F_d$-proofs. Then for every $m\in M\setminus\mathbb N$, there exists a $\chi\subseteq\mathcal P(M_m)$ such that:
\begin{enumerate}
\item Every bounded subset of $M_m$ which is definable in $M$ is in $\chi$.
\item $(M_m,\chi)\models V_1^0$.
\item There is an $\alpha\in\chi$ such that $(M_m,\chi)\models\neg\varphi(m,\alpha)$.
\end{enumerate}
\section*{3 Razborov's conjecture for $\mathrm{AC}^0$-Frege}
In this section, we state and prove the main result of the paper.
A Boolean function $f$ is symmetric iff for every $n$, $f_n$ is invariant under any permutation of its inputs. Let $S_{\mathrm{AC}_2^0}$ denote the depth two $\mathrm{AC}^0$ circuit complexity of functions, then:
\paragraph{Theorem 6.} Let $f\in\mathrm{NP}\cap\mathrm{CoNP}$ be a symmetric function such that $S_{\mathrm{AC}_2^0}(f)=2^{n^{\Omega(1)}}$ and $(\varphi_0,\varphi_1)$ be a $\Sigma_1^1\cap\Pi_1^1$ definition of $f$. Then for every $d$:
\begin{enumerate}
\item For every positive $c\in\mathbb N$, every $0<\epsilon<1$, there exits an $\epsilon'>0$ such that for every large enough $n$, every $n^c\times n$ $\lfloor n^\epsilon\rfloor$-sparse matrix $A$, any $b\notin\operatorname{Rng}$, $\tau_b(\mathrm{NW}_{f,A})$ based on $(\langle\varphi_0\rangle_{\lfloor n^\epsilon\rfloor},\langle\varphi_1\rangle_{\lfloor n^\epsilon\rfloor})$ does not have $F_d$-proofs of size less than $2^{n^{\epsilon'}}$.
\item For every positive $r\in\mathbb N$, every large enough $s$, every $t\in[s/r,s]$, every $c\in\mathbb N$, every large enough $n$, every $2^n\times n^s$ $n^t$-sparse matrix $A$, any $b\notin\operatorname{Rng}$, $\tau_b(\mathrm{NW}_{f,A})$ based on $(\langle\varphi_0\rangle_{n^t},\langle\varphi_1\rangle_{n^t})$ does not have $F_d$-proofs of size less than $2^{cn}$.
\end{enumerate}
Note that in Theorem 6, the size of the formula $\tau_b(\mathrm{NW}_{f,A})$ is $n^{O(1)}$ in the first part and it is $2^{O(n)}$ in the second part. This theorem is proved by a model-theoretic argument based on the relations explained in Preliminaries in combination with the hardness of the Pigeonhole principle in $\mathrm{AC}^0$-Frege. Model theoretic arguments have been used previously in proof complexity and they were very fruitful (for example see [1,11,17] and [9] for a detailed explanation). See [39,12] for discussions about the importance and benefits of the model-theoretic arguments (and in general, the logical point of view) in proof complexity.
Note that an immediate consequence of Theorem 6 is that $\mathrm{NW}_{f,A}$ based on a hard enough function $f$ with suitable parameters is a hard proof complexity generator for $\mathrm{AC}^0$-Frege. As the Parity function or the Majority function satisfies the required assumptions of Theorem 6, we get that NW-generators based on these functions are hard proof complexity generators for $\mathrm{AC}^0$-Frege.
\paragraph{Proof of Theorem 6.} We state the proof as a series of lemmas for more clarity. We prove the second part of this theorem. The first part can be proved in the same way. For the rest of the paper, $[n]:=[0,n)$.
Intuitively, the proof goes as follows. Let $n$ be large enough. As $f$ is symmetric and has high depth two circuit complexity, then there is a $u$ such that $f_{n^t}$ on strings with $u$ many $1$'s is different from $f_{n^t}$ on strings with $u+1$ many $1$'s and moreover, $u$ is not too big and not too small. This implies that there is an input $a$ for the NW generator such that the number of $1$'s of $a$ restricted to each row is close to $u$. Now, assuming PHP fails, for every row $i$ we can find a witness for $\varphi_{b_i}(a|_{J_i(A)})$ by constructing permutations between sets that actually have different sizes and this forces any string $b$ into the range of the generator and this is possible as the truth of representations are closed under permutations. The next lemmas formalize this intuitive proof and their combination proves the theorem.
\paragraph{Lemma 7.} Let $f:\{0,1\}^*\to\{0,1\}$ be a symmetric Boolean function such that $S_{\mathrm{AC}_2^0}(f_n)=\Omega(2^{n^\epsilon})$ for an $\epsilon>0$. Then there is a natural $m$ such that for every $n\geq m$ there is natural number $u\in[n^{\epsilon/2},n-n^{\epsilon/2}]$ such that
\[f_n(1^u0^{n-u})\neq f_n(1^{u+1}0^{n-u-1}).\]
\paragraph{Proof.} Let $g:\{0,1\}^n\to\{0,1\}$ be a symmetric function. If there exists a $r\leq n/2$ such that for every $r\leq k\leq n-r$, $g(1^k0^{n-k})=0$, then
\[S_{\mathrm{DNF}}(g)\leq2n\cdot\sum_{i=0}^r\binom ni\leq2n(en/r)^r\]
where $S_{\mathrm{DNF}}$ denotes the DNF complexity of functions. Writing this inequality (1) for $f_n$, we get $c2^{n^\epsilon}\leq2n(en/r)^r$ for a $c>0$. So if we put $r=n^\delta$ and rewriting this inequality we have
\[c2^{n^\epsilon}\leq2n(en^{1-\delta})^{n^\delta}\leq2e^{n^\delta}n^{1+n^\delta}\leq2n^{1+2n^\delta}=2^{(2n^\delta+1)\log n+1}.\]
So assuming $\delta=\epsilon/2$, we have $(2n^\delta+1)\log n+1=o(n^\epsilon)$. Therefore for every large enough $n$, there exists a $v\in[n^{\epsilon/2},n-n^{\epsilon/2}]$ such that $f_n(1^v0^{n-v})=1$. Following the same argument for $\neg f_n$, we can deduce that for every large enough $n$, there exists a $v'\in[n^{\epsilon/2},n-n^{\epsilon/2}]$ such that $\neg f_n(1^{v'}0^{n-v'})=1$. So we have found $v,v'\in[n^{\epsilon/2},n-n^{\epsilon/2}]$ such that $f_n(1^v0^{n-v})\neq f_n(1^{v'}0^{n-v'})$, hence there exists a $u\in[n^{\epsilon/2},n-n^{\epsilon/2}]$ such that
\[f_n(1^u0^{n-u})\neq f_n(1^{u+1}0^{n-u-1}).\]\hfill$\triangleleft$
Now let $M$ be a countable nonstandard model of true arithmetic. Let $n,s,t,A,b$ be arbitrary elements of $M$ such that:
\begin{enumerate}
\item $n,t\in M\setminus\mathbb N$.
\item $A\in M\setminus\mathbb N$ encodes a $2^n\times n^s$ $n^t$-sparse matrix where $t\in[s/r,s]$, $n^s<2^n$, and $n^s2^n\leq2^{n^{t/u}}$ for a nonstandard $u$.
\item $b\in M\setminus\mathbb N$ is a binary string of length $2^n$ such that $b\notin\operatorname{Rng}(\mathrm{NW}_{f,A})$.
\end{enumerate}
Let $\chi$ be the set of all bounded subsets of $M_{n^t}$ encoded in $M$. So in particular $A,b\in\chi$.
As $S_{\mathrm{AC}_2^0}(f_m)=2^{m^{\Omega(1)}}$, there is a standard rational $\epsilon>0$ such that $S_{\mathrm{AC}_2^0}(f_m)=\Omega(2^{m^\epsilon})$. Let $\delta:=\epsilon/2$, then there exists $u\in[n^{\delta t},n^t-n^{\delta t}]$ that is guaranteed to exist by Lemma 7 for $f_{n^t}$. Let $v:=\min\{u,n^t-u\}$, then
\paragraph{Lemma 8.} There exists a binary string $\alpha\in\chi$ of length $n^s$ such that for every $i\in[2^n]$,
\[\#_1(\alpha|_{J_i(A)})\in[v(1-1/\sqrt[3]v),v(1+1/\sqrt[3]v)]\]
where $\#_s(w)$ is the number of occurrences of symbol $s$ in the string $w$.
\paragraph{Proof.} Let $X_0,\ldots,X_{n^s-1}$ be independent random variables taking values in $\{0,1\}$ such that for every $i$, $\Pr[X_i=1]=v/n^t$. For every $i\in[2^n]$, let $Y_i=\sum_{j\in J_i(A)}X_j$ and hence $\mathbb E[Y_i]=v$.
By the Chernoff bound we have the following inequalities for every $i\in[2^n]$:
\begin{enumerate}
\item $\Pr[Y_i\leq v(1-1/\sqrt[3]v)]\leq e^{-\sqrt[3]v/2}$.
\item $\Pr[Y_i\geq v(1+1/\sqrt[3]v)]\leq e^{-\sqrt[3]v/3}$.
\end{enumerate}
Let $X'$ be the concatenation of $X_0,\ldots,X_{n^s-1}$, hence it is a random string of length of $n^s$. Now combining the above inequalities with the union bound we get:
\begin{align*}
P&=\Pr\left[\bigvee_{i=0}^{2^n-1}\#_1(X'|_{J_i(A)})\notin[v(1-1/\sqrt[3]v),v(1+1/\sqrt[3]v)]\right]\\
&\leq\sum_{i=0}^{2^n-1}\Pr[\#_1(X'|_{J_i(A)})\notin[v(1-1/\sqrt[3]v),v(1+1/\sqrt[3]v)]]\\
&\leq\sum_{i=0}^{2^n-1}(\Pr[Y_i\leq v(1-1/\sqrt[3]v)]+\Pr[Y_i\geq v(1+1/\sqrt[3]v)])\\
&\leq2^n\cdot2e^{-\sqrt[3]v/3}.
\end{align*}
We know that $v\geq n^{\delta t}$, $t$ is a nonstandard number, and $\delta$ is a standard rational, so
\[n+1<n^{\delta t/3}/3\leq\sqrt[3]v/3\]
which implies $2^n\cdot2e^{-\sqrt[3]v/3}<1$, and hence $P<1$. This implies that there exists a string $\alpha\in\chi$ that satisfies the desired property.\hfill$\triangleleft$
\paragraph{Lemma 9.} The following functions exist in $\chi$:
\begin{enumerate}
\item $\gamma:[2^n]\to[n^t+1]$ such that for every $i\in[2^n]$, $\gamma(i)=\#_1(\alpha|_{J_i(A)})$.
\item $\omega:[2^n]\times[n^t]\to[n^t]$ such that for every $i\in[2^n]$, $\omega(i,\cdot)$ defines a permutation over $[n^t]$ and moreover $\beta_j=\alpha|_{J_i(A)}(\omega(i,j))$ where $\beta=1^{\gamma(i)}0^{n^t-\gamma(i)}$.
\end{enumerate}
\paragraph{Proof.}
\begin{enumerate}
\item The function $\gamma$ exists in $M$ as it is definable by an arithmetical formula with parameter $\alpha$. To prove that $\gamma$ is in $\chi$, we observe that encoding of $\gamma$ as a binary string requires at most $c2^n\cdot\log n^t$ (for some $c\in\mathbb N$) which is less than $2^{n^{\sqrt t}}$, hence $\gamma\in\chi$.
\item Like the previous part, $\omega$ exists in $M$ as it is definable by an arithmetical formula with parameters $\alpha$ and $\gamma$, and its bit representation requires at most $c2^n\cdot n^t\log n^t$ (for some $c\in\mathbb N$) which is again less than $2^{n^{\sqrt t}}$, and therefore $\omega\in\chi$.
\end{enumerate}\hfill$\triangleleft$
To continue the proof, we need the celebrated result about the hardness of the Pigeonhole principle for $\mathrm{AC}^0$-Frege.
\paragraph{Theorem 10 ([1,25,33]).} For any natural number $d$, there exists an $\epsilon_d>0$ such that for large values of $n$, any $F_d$-proof of $\mathrm{PHP}_n^{n+1}$ has size at least $2^{\Omega(n^{\epsilon_d})}$.

Now let $l=\lfloor\sqrt[4]v\rfloor$, then we have the following lemma.
\paragraph{Lemma 11.} There exists $\chi'\subseteq\mathcal P(M_{n^t})$ such that:
\begin{enumerate}
\item $\chi\subseteq\chi'$.
\item There exists a function $\sigma\in\chi'$ such that $\sigma$ is a bijection from $[l]$ onto $[l-1]$.
\item $(M_{n^t},\chi')\models V_1^0$.
\end{enumerate}
\paragraph{Proof.} By Theorem 10 we know that $\mathrm{PHP}_{m-1}^m$ requires exponential size $F_d$-proofs for every $d$. Therefore by Theorem 5, there exists a $\chi'\subseteq\mathcal P(M_l)$ such that every bounded subset of $M_l$ is in $\chi'$, $(M_l,\chi')\models V_1^0$ and there exists a $\sigma\in\chi'$ such that it is bijection from $[l]$ onto $[l-1]$.
Note that if $a\in M_{n^t}$, then there exists a $b\in M\setminus\mathbb N$ such that $a<2^{n^{t/b}}$. Let $b'=\lfloor\delta b/4\rfloor$, then $a<2^{l^{1/b'}}$ as we know $l\geq n^{\delta t/4}$. This implies that $M_l=M_{n^t}$, and moreover $\chi\subseteq\chi'$ which completes the proof.\hfill$\triangleleft$
The following lemma shows that we can simultaneously falsify some weak Pigeonhole principle instances.
\paragraph{Lemma 12.} There exists a function $F:[2^n]\times\{-1,0,1\}\times[n^t]\to[n^t]$ in $\chi'$ such that for every $i\in[2^n]$
\begin{enumerate}
\item $F(i,a,\cdot)$ restricted to $[v+a]$, is a bijection from $[v+a]$ onto $[\gamma(i)]$.
\item $F(i,a,\cdot)$ restricted to $[v+a,n^t)$, is a bijection from $[v+a,n^t)$ onto $[\gamma(i),n^t)$.
\end{enumerate}
\paragraph{Proof.} Let $\sigma\in\chi'$ be the function that Lemma 11 provides. Let
\begin{enumerate}
\item $w_{i,a}=|\gamma(i)-v-a|$.
\item $M_{i,a}=\max\{v+a,\gamma(i)\}$.
\item $m_{i,a}=\min\{v+a,\gamma(i)\}$.
\end{enumerate}
Then we define the function $G_0(i,a,b)$ as follows:
\[G_0(i,a,b):=\begin{cases}\sigma(b-lk)+(l-1)k&b\in[lk,l(k+1))\wedge k\in[w_{i,a}]\\b-w_{i,a}&b\in[lw_{i,a},M_{i,a})\end{cases}\]
where $i\in[2^n]$, $a\in\{-1,0,1\}$, and $b\in[M_{i,a}]$.
Note that $v-1\leq M_{i,a}$, hence
\[w_{i,a}\leq\frac{v-1}{\sqrt[4]v}\leq\frac{v-1}l\leq\frac{M_{i,a}}l\]
as $w_{i,a}\leq\sqrt[3]{v^2}+1$ by Lemma 8. So $G_0(i,a,\cdot)$ is a bijection from $[lw_{i,a}]$ onto $[(l-1)w_{i,a}]$ and moreover is a bijection from $[lw_{i,a},M_{i,a})$ onto $[(l-1)w_{i,a},m_{i,a})$ as
\[M_{i,a}-lw_{i,a}=m_{i,a}-(l-1)w_{i,a}.\]
Therefore the conclusion is that $G(i,a,\cdot)$ is a bijection from $[M_{i,a}]$ onto $[m_{i,a}]$ for every $i\in[2^n]$ and $a\in\{-1,0,1\}$. Now, we define the function $G_1(i,a,b)$ as the inverse of $G_0$ which means that
\[G_1(i,a,G_0(i,a,b))=b\]
where $i\in[2^n]$, $a\in\{-1,0,1\}$, and $b\in[M_{i,a}]$. So $G_1(i,a,\cdot)$ is a bijection from $[m_{i,a}]$ onto $[M_{i,a}]$.
Using $G_0$ and $G_1$, we can fulfill (1) from the lemma. Now we want to construct two other functions $H_0$ and $H_1$ to fulfill (2).
The task is to define $H_0(i,a,\cdot)$ as a function that defines a bijection from $[\max\{n^t-v-a,n^t-\gamma(i)\}]$ onto $[\min\{n^t-v-a,n^t-\gamma(i)\}]$ where $i\in[2^n]$ and $a\in\{-1,0,1\}$ and moreover $H_1$ would be the inverse of $H_0$. Let
\begin{enumerate}
\item $M'_{i,a}=\max\{n^t-v-a,n^t-\gamma(i)\}$.
\item $m'_{i,a}=\min\{n^t-v-a,n^t-\gamma(i)\}$.
\end{enumerate}
Then we define $H_0(i,a,b)$ as follows:
\[H_0(i,a,b):=\begin{cases}\sigma(b-lk)+(l-1)k&b\in[lk,l(k+1))\wedge k\in[w_{i,a}]\\b-w_{i,a}&b\in[lw_{i,a},M'_{i,a})\end{cases}\]
where $i\in[2^n]$, $a\in\{-1,0,1\}$, and $b\in[M'_{i,a}]$.
Note that $n^t-v-1\leq M'_{i,a}$ and moreover $v\leq n^t/2$, therefore $n^t/2-1\leq M'_{i,a}$. This implies that
\begin{align*}
w_{i,a}&\leq\sqrt[3]{v^2}+1\leq\sqrt[3]{(n^t/2)^2}+1\leq\sqrt[4]{(n^t/2)^3}-1\\
&\leq\frac{n^t/2-1}{\sqrt[4]{n^t/2}}\leq\frac{n^t/2-1}{\sqrt[4]v}\leq\frac{n^t/2-1}l\leq\frac{M'_{i,a}}l.
\end{align*}
Therefore $H_0(i,a,\cdot)$ is a bijection from $[lw_{i,a}]$ onto $[(l-1)w_{i,a}]$ and moreover is a bijection $[lw_{i,a},M'_{i,a})$ onto $[(l-1)w_{i,a},m'_{i,a})$. Hence $H_0(i,a,\cdot)$ is a bijection from $[M'_{i,a}]$ onto $[m'_{i,a}]$ for every $i\in[2^n]$ and $a\in\{-1,0,1\}$. Now we define the function $H_1(i,a,b)$ as the inverse again as follows:
\[H_1(i,a,H_0(i,a,b))=b\]
where $i\in[2^n]$, $a\in\{-1,0,1\}$, and $b\in[M'_{i,a}]$. Hence $H_1(i,a,\cdot)$ is a bijection from $[m'_{i,a}]$ onto $[M'_{i,a}]$.
Now $F(i,a,b)$ is:
\[F(i,a,b):=\begin{cases}G_0(i,a,b)&b\in[v+a]\wedge v+a=M_{i,a}\\G_1(i,a,b)&b\in[v+a]\wedge v+a=m_{i,a}\\H_0(i,a,b-v-a)&b\in[v+a,n^t)\wedge n^t-v-a=M'_{i,a}\\H_1(i,a,b-v-a)&b\in[v+a,n^t)\wedge n^t-v-a=m'_{i,a}.\end{cases}\]
As $G_0,G_1,H_0,H_1$ are definable by a bounded arithmetical formula based on $\gamma$ and $\sigma$, therefore $F$ is also definable by a bounded arithmetical formula based on $\gamma$ and $\sigma$ and this implies that $F\in\chi'$ as $(M_{n^t},\chi')\models V_1^0$.\hfill$\triangleleft$
Without the loss of generality we can assume $f_{n^t}(1^u0^{n^t-u})=0$. Consider the following relations in $\chi$:
\begin{enumerate}
\item $\theta_0=1^u0^{n^t-u}$.
\item $\theta'_0=0^{n^t-u}1^u$.
\item $\theta_1=1^{u+1}0^{n^t-u-1}$.
\item $\theta'_1=0^{n^t-u-1}1^{u+1}$.
\item $\overline\lambda_0$ such that $\varphi_0(\theta_0,\overline\lambda_0)$ holds in $(M_{n^t},\chi)$.
\item $\overline\lambda'_0$ such that $\varphi_0(\theta'_0,\overline\lambda'_0)$ holds in $(M_{n^t},\chi)$.
\item $\overline\lambda_1$ such that $\varphi_1(\theta_1,\overline\lambda_1)$ holds in $(M_{n^t},\chi)$.
\item $\overline\lambda'_1$ such that $\varphi_1(\theta'_1,\overline\lambda'_1)$ holds in $(M_{n^t},\chi)$.
\end{enumerate}
Now we are ready to describe the assignments $\mathcal X$, $\{\overline Y_i\}_{i\in[2^n]}$, and $\{\overline Z_i\}_{i\in[2^n]}$ such that
\[(M_{n^t},\chi')\models\forall i<2^n\ (b_i=0\to\varphi_0(\mathcal X|_{J_i(A)},\overline Y_i))\wedge(b_i=1\to\varphi_1(\mathcal X|_{J_i(A)},\overline Z_i))\]
which implies that $\tau_b(\mathrm{NW}_{f,A})$ based on $(\langle\varphi_0\rangle_{n^t},\langle\varphi_1\rangle_{n^t})$ fails under an assignment in $(M_{n^t},\chi')$. We define these assignments as follows:
\begin{enumerate}
\item If $u=v$:
\begin{enumerate}
\item $\mathcal X=\alpha$.
\item $\overline Y_i=\omega(i,F(i,0,\overline\lambda_0))$.
\item $\overline Z_i=\omega(i,F(i,1,\overline\lambda_1))$.
\end{enumerate}
\item If $u=n^t-v$:
\begin{enumerate}
\item $\mathcal X$ is the complement of $\alpha$, i.e., $\mathcal X_j=1-\alpha_j\ j\in[n^s]$.
\item $\overline Y_i=\omega(i,F(i,0,\overline\lambda'_0))$.
\item $\overline Z_i=\omega(i,F(i,-1,\overline\lambda'_1))$.
\end{enumerate}
\end{enumerate}
Without loss of generality assume $v=u$. Then for an arbitrary $i\in[2^n]$, we know that
\[\mathcal X|_{J_i(A)}=\omega(i,F(i,0,\theta_0)),\]
hence $\sigma_0(\mathcal X|_{J_i(A)},\overline Y_i)$ holds by Lemma 3 as $\omega(i,F(i,0,\cdot))$ is a bijection from $[n^t]$ onto itself and the fact that $\sigma_0(\theta_0,\overline\lambda_0)$ holds. The same argument works for $\sigma_1(\mathcal X|_{J_i(A)},\overline Z_i)$. Moreover if $v=n^t-u$, the same argument works by using $\theta'_0,\theta'_1,\overline\lambda'_0,\overline\lambda'_1$.
To complete the proof, we argue as follows. Suppose the statement of the theorem is not true. This means that there exist standard $d$ and $r$ such that the following arithmetical sentence is true in $\mathbb N$:
\begin{multline*}
H:=\forall s_1\exists s\geq s_1,\exists t\in[s/r,s],\exists c>0,\forall m,\exists n>m\exists\ 2^n\times n^s\ n^t\text{-sparse matrix }A,\\
\exists b\notin\operatorname{Rng}(\mathrm{NW}_{f,A}),\exists F_d\text{-proof }\pi\text{ for }\tau_b(\mathrm{NW}_{f,A})\text{ such that }|\pi|\leq|\tau_b(\mathrm{NW}_{f,A})|^c.
\end{multline*}
Let $M$ be a countable nonstandard model of true arithmetic. This means that $M\models H$. To simplify the presentation let
\[H:=\forall s_1\exists s,t,c\forall m\exists n,A,b,\pi\ \Phi(s_1,s,t,c,m,n,A,b,\pi).\]
Let $s_1\in M\setminus\mathbb N$, then there exist $s,t\in M\setminus\mathbb N$ and $c\in M$ such that
\[M\models\forall m\exists n,A,b,\pi\ \Phi(s_1,s,t,c,m,n,A,b,\pi).\]
We choose an $m\in M\setminus\mathbb N$ such that for all $m_1\geq m$, $m_1^{ct}2^{cm_1}\leq2^{m_1^{\sqrt{t/2}}}$, hence there exist an $n>m$, a $2^n\times n^s$ $n^t$-sparse matrix $A\in M\setminus\mathbb N$, a $b\in M\setminus\mathbb N$ such that $b\notin\operatorname{Rng}(\mathrm{NW}_{f,A})$, and an $F_d$-proof $\pi\in M$ for $\tau_b(\mathrm{NW}_{f,A})$ such that $|\pi|\leq|\tau_b(\mathrm{NW}_{f,A})|^c$. Now we consider $M_{n^t}$ and by the argument in this section, there exists a $\chi'\subseteq\mathcal P(M_{n^t})$ such that it has every bounded $M$-definable subset of $M_{n^t}$ and moreover
\begin{enumerate}
\item $(M_{n^t},\chi')\models V_1^0$.
\item There exists an $\alpha\in\chi'$ which falsifies $\tau_b(\mathrm{NW}_{f,A})$.
\end{enumerate}
Then by Theorem 4 there is no $F_d$-proof of $\tau_b(\mathrm{NW}_{f,A})$ in $(M_{n^t},\chi')$. Note that there is a standard number $e$ such that
\[|\pi|\leq|\tau_b(\mathrm{NW}_{f,A})|^c\leq(n^{ct}2^{cn})^e\leq2^{en^{\sqrt{t/2}}}\]
which implies that $\pi\in\chi$, but this leads to a contradiction and completes the proof.
